\subsection{Dhclient Script}
\begin{figure}[htp]
\centering
\includegraphics[width=0.5\textwidth]{dhclient_service_script}
\label{fig:dhclient_script_statements}
\caption{Dhclient Script Statements}
\end{figure}


\TheStatement{dhclient}
\TheStatement*[dhclient]{dhclient \{ requests prepend \}}

Entry point in the dhclient script.

%% requests
\TheStatement{requests}
\TheStatement*[dhclient!requests]{requests \Arg{requests}}

Sets the list of DHCP \Arg{requests}. Request item that starts with an 
exclamation mark \qcode{!} is removed from the list. The default list
of requests is: 
\begin{inparaitem}[]
\item subnet-mask,
\item broadcast-address,
\item time-offset,
\item routers,
\item domain-name,
\item domain-name-servers,
\item domain-search,
\item host-name,
\item netbios-name-servers,
\item netbios-scope,
\item interface-mtu,
\item rfc3442-classless-static-routes,
\item ntp-servers.
\end{inparaitem}

\TheStatement[prepend]{prepend}
\TheStatement*[dhclient!prepend]{prepend \Arg{option} \Arg{desclaration}}

Adds the \Arg{option} with the \Arg{desclaration} as a prepend statement 
to the DHCP configuration.

